\documentclass{article}
\usepackage[T1]{fontenc}
\usepackage{lmodern}
\usepackage{graphicx}
\usepackage{amsmath,amssymb}
\usepackage[a4paper,margin=2cm]{geometry}
\usepackage[absolute,overlay]{textpos}
\setlength{\TPHorizModule}{1mm}
\setlength{\TPVertModule}{1mm}
\usepackage[indonesian]{babel}
\usepackage{hyperref}
\setlength{\emergencystretch}{2em}
% unit: Halaman 56

\begin{document}

% === Halaman 56 (python/native) ===

Contoh 1:  Misalkan a = 00001101 = x3 + x2 + 1 dan b = 00000110 = x2 + x 


       a dan b  GF(28)      (dalam notasi Hex, a = ’0D’ dan b = ’06’)

(i)

a + b = ‘0D’ + ‘06’ =  (x3 + x2 + 1) +  (x2 + x ) = x3 + 2x2 + x + 1 (mod 2)







   = x3 + 0x2 + x + 1 







   = x3 + x + 1  = 00001011 = ‘0B’


Dalam representasi biner:


00001101


00000110 +


00001011  → sama dengan hasil operasi XOR


 a + b = 00001011 = a XOR b 

Bagi  tiap koefisien dengan 2,

lalu ambil sisanya

\end{document}
